% =====================================================================
%  Kreyszig section notes: BLANK template
%  1. Copy and rename, e.g. 11-2-arbitrary-period.tex
%  2. Fill in the 6 lines under "Fill in"
%  3. Write under each heading; delete any section you don't need
% =====================================================================
\documentclass[11pt,a4paper]{article}

\usepackage[margin=2.2cm]{geometry}
\usepackage{amsmath,amssymb,amsthm}
\usepackage{enumitem}
\usepackage[most]{tcolorbox}
\usepackage{fancyhdr}
\usepackage{hyperref}

% ---------- Fill in --------------------------------------------------
\newcommand{\SectionNo}{11.2}
\newcommand{\SectionTitle}{Fourier Series: Arbitrary Period}
\newcommand{\Lesson}{01.20}
\newcommand{\Sprint}{S1}
\newcommand{\DateRead}{2026-11-12}
\newcommand{\Confidence}{}       % e.g. 3/5

% ---------- Shortcuts ------------------------------------------------
\newcommand{\vect}[1]{\mathbf{#1}}                   % \vect{a}
\newcommand{\pd}[2]{\frac{\partial #1}{\partial #2}} % \pd{f}{x}
\newcommand{\dv}[2]{\frac{\mathrm{d} #1}{\mathrm{d} #2}} % \dv{y}{x}
\newcommand{\dd}{\,\mathrm{d}}                       % \int f \dd x
\newcommand{\R}{\mathbb{R}}
\newcommand{\C}{\mathbb{C}}

% ---------- Boxes ----------------------------------------------------
\tcbset{boxrule=0.6pt, arc=2pt, left=6pt, right=6pt, top=4pt, bottom=4pt, fonttitle=\bfseries}
\newtcolorbox{formulabox}[1][Key formulas]{colback=green!4, colframe=green!40!black, title=#1}
\newtcolorbox{mistakebox}{colback=red!4, colframe=red!50!black, title=Mistakes and fixes}
\newtcolorbox{mlbox}{colback=orange!5, colframe=orange!60!black, title=ML / robotics link}
\newtcolorbox{keybox}{colback=blue!4, colframe=blue!50!black, title=If I remember only one thing}

\theoremstyle{definition}
\newtheorem{definition}{Definition}
\newtheorem{example}{Example}
\theoremstyle{plain}
\newtheorem{theorem}{Theorem}

\pagestyle{fancy}
\fancyhf{}
\lhead{Kreyszig \SectionNo}
\rhead{Lesson \Lesson\ \textbullet\ \Sprint}
\cfoot{\thepage}
\setlist{itemsep=2pt, topsep=4pt}

% ---------- Title --------------------------------------------------
\title{\SectionNo\quad \SectionTitle}
\author{Ravikant Pothukuchi}          % your name here
\date{\DateRead}

\begin{document}
\maketitle


% ---------------------------------------------------------------------
\section{Introduction}
% 1-2 sentences in your own words.
Arbitrary period. For Fourier series with period $2\pi$ it was easy. 
We build a series and find coefficients.
These are called Euler formulas probably because he discovered them.

\begin{formulabox}[Fourier series and Euler formulas (period $2\pi$)]
\begin{align}
  f(x) &= a_0 + \sum_{n=1}^{\infty} \bigl(a_n \cos nx + b_n \sin nx\bigr)\\
  a_0 &= \frac{1}{2\pi}\int_{-\pi}^{\pi} f(x)\dd x \qquad\text{(average value of } f)\\
  a_n &= \frac{1}{\pi}\int_{-\pi}^{\pi} f(x)\cos nx\dd x, \qquad
  b_n = \frac{1}{\pi}\int_{-\pi}^{\pi} f(x)\sin nx\dd x
\end{align}
\end{formulabox}


Now we enter domain where periods can be $2L$. 
\begin{formulabox}[Fourier series and Euler formulas (period $2L$)]
\begin{align*}
  f(x) &= a_0 + \sum_{n=1}^{\infty}\left(a_n \cos\frac{n\pi x}{L} + b_n \sin\frac{n\pi x}{L}\right)\\[6pt]
  a_0 &= \frac{1}{2L}\int_{-L}^{L} f(x)\dd x \tag{0}\\
  a_n &= \frac{1}{L}\int_{-L}^{L} f(x)\cos\frac{n\pi x}{L}\dd x && n = 1, 2, \dots \tag{a}\\
  b_n &= \frac{1}{L}\int_{-L}^{L} f(x)\sin\frac{n\pi x}{L}\dd x && n = 1, 2, \dots \tag{b}
\end{align*}
\end{formulabox}

% ---------------------------------------------------------------------
\section{The Twist}
% \begin{definition}[Name] ... \end{definition}
The way this was derived is not important. Its a long convulated way of using $2\pi$ to get a different variable v and rebuilding the original formula. Its a side quest if you want to pursue.



\end{document}